\documentclass[../main]{subfiles}
\begin{document}

\chapter{Motor Asincrónico Trifásico}

\section{Motor Asincrónico Trifásico}

    \subsection{Introducción al Motor Asincrónico Trifásico}

    Un motor asincrónico trifásico, también conocido como motor de inducción, es una máquina eléctrica que convierte energía eléctrica en energía mecánica mediante la interacción entre campos magnéticos. Es ampliamente utilizado en aplicaciones industriales debido a su robustez, simplicidad y eficiencia. Este tipo de motor opera sin necesidad de una conexión eléctrica directa entre el rotor y el estator, lo que reduce el desgaste y la necesidad de mantenimiento.

    \info{Principio de Funcionamiento}{
    El motor asincrónico trifásico funciona basado en el principio de inducción electromagnética. Cuando se aplica un sistema de tensiones trifásicas al estator, se genera un campo magnético giratorio. Este campo induce una corriente en el rotor, creando un segundo campo magnético. La interacción entre estos campos produce un par motor que hace girar el rotor.

    La velocidad del campo magnético giratorio, llamada velocidad síncrona ($n_s$), se calcula mediante la fórmula:
    \begin{equation}
    n_s = \frac{120 \times f}{P}
    \end{equation}
    Donde:
    \begin{itemize}
        \item $n_s$: Velocidad síncrona en revoluciones por minuto (RPM).
        \item $f$: Frecuencia de la fuente de alimentación en Hertz (Hz).
        \item $P$: Número de polos del motor.
    \end{itemize}
    }

    \info{Deslizamiento}{
    El deslizamiento ($s$) es una medida de la diferencia entre la velocidad síncrona y la velocidad real del rotor ($n_r$). Se expresa como un porcentaje y se calcula con la siguiente fórmula:
    \begin{equation}
    s = \frac{n_s - n_r}{n_s} \times 100
    \end{equation}
    Donde:
    \begin{itemize}
        \item $s$: Deslizamiento en porcentaje.
        \item $n_s$: Velocidad síncrona (RPM).
        \item $n_r$: Velocidad del rotor (RPM).
    \end{itemize}
    }

    \info{Potencia del Motor}{
    La potencia mecánica ($P_m$) desarrollada por el motor se puede calcular usando la fórmula:
    \begin{equation}
    P_m = T \times \omega
    \end{equation}
    Donde:
    \begin{itemize}
        \item $P_m$: Potencia mecánica en vatios (W).
        \item $T$: Par motor en Newton-metro (Nm).
        \item $\omega$: Velocidad angular en radianes por segundo (rad/s).
    \end{itemize}
    }

\actividad{Responder el siguiente cuestionario}
\begin{enumerate}
    \item ¿Qué es un motor asincrónico trifásico y cuáles son sus principales aplicaciones?
    \item ¿Cómo se genera el campo magnético giratorio en un motor asincrónico trifásico?
    \item ¿Qué es la velocidad síncrona y cómo se calcula?
    \item Explica el concepto de deslizamiento en un motor asincrónico.
    \item ¿Por qué es importante el deslizamiento en el funcionamiento de un motor asincrónico?
    \item ¿Qué relación hay entre el par motor y la potencia mecánica en un motor asincrónico?
    \item ¿Qué ventajas tiene un motor asincrónico trifásico sobre otros tipos de motores eléctricos?
    \item ¿Qué factores afectan la eficiencia de un motor asincrónico trifásico?
    \item ¿Cómo se puede controlar la velocidad de un motor asincrónico trifásico?
    \item ¿Qué tipos de mantenimiento se recomiendan para un motor asincrónico trifásico?

\end{enumerate}

\actividad{Resolver los siguientes problemas}
\begin{enumerate}
    \item Un motor asincrónico trifásico tiene una frecuencia de 50 Hz y 4 polos. ¿Cuál es su velocidad síncrona?
    \item Un motor con una velocidad síncrona de 1500 RPM tiene un deslizamiento del 3\%. ¿Cuál es la velocidad del rotor?
    \item Si un motor desarrolla un par motor de 25 Nm a una velocidad de 1450 RPM, ¿cuál es la potencia mecánica desarrollada?
    \item Un motor de 6 polos está conectado a una fuente de 60 Hz. ¿Cuál es la velocidad síncrona del motor?
    \item La velocidad del rotor de un motor es de 1425 RPM y su velocidad síncrona es de 1500 RPM. ¿Cuál es el deslizamiento?
    \item Un motor asincrónico trifásico tiene una velocidad del rotor de 1440 RPM y un deslizamiento del 4\%. ¿Cuál es su velocidad síncrona?
    \item Calcula la potencia mecánica desarrollada por un motor que tiene un par motor de 30 Nm y una velocidad de 1400 RPM.
    \item Un motor con una velocidad síncrona de 1800 RPM tiene un deslizamiento del 2\%. ¿Cuál es la velocidad del rotor?
    \item Si un motor desarrolla una potencia mecánica de 1500 W con una velocidad angular de 157 rad/s, ¿cuál es el par motor?
    \item Un motor asincrónico trifásico tiene un deslizamiento del 5\% y una velocidad del rotor de 1430 RPM. ¿Cuál es su velocidad síncrona?

\end{enumerate}

\end{document}
